problem:  
Let \((M^n, g)\) be a compact, simply connected **Einstein manifold** of dimension \(n \ge 5\) with **positive sectional curvature**, and let  
\[
R = W + U + V
\]
denote the standard curvature decomposition into Weyl, traceless Ricci, and scalar components.  
In the Einstein case, \(U = 0\), so \(R = W + V\).  

Assume that the Weyl curvature satisfies a smallness condition in the **critical conformally invariant norm**:
\[
\|W\|_{L^{n/2}(M)} \le \varepsilon(n) \, \|R\|_{L^{n/2}(M)},
\]
where \(\varepsilon(n) > 0\) is a dimensional constant.  

**Question (Weyl–Energy Gap Conjecture for Positively Curved Einstein Metrics):**  
Is there a universal constant \(\varepsilon(n)\) such that whenever the above inequality holds, \((M,g)\) must be **locally isometric to a compact rank–one symmetric space (CROSS)**, i.e., the round sphere \(S^n\), complex projective space \(\mathbb{CP}^{n/2}\), quaternionic projective space \(\mathbb{HP}^{n/4}\), or the Cayley plane \(\mathrm{CaP}^2\)?  

Why is it a "good" problem:  
This formulation establishes a precise **Weyl–energy rigidity conjecture** that operates at the intersection of conformal geometry, global analysis, and Riemannian topology. It seeks to extend known “pinching implies rigidity” phenomena into a setting governed by the scale–invariant \(L^{n/2}\) geometry of the Weyl tensor.

The problem integrates several deep ideas:
1. **Analytic–geometric bridge:** It proposes that small conformal distortion (measured by the Weyl tensor in the critical norm) could enforce **global symmetry** under positive curvature and the Einstein constraint—an analytic analogue of the classical **sphere theorem**, yet formulated in a **conformal–integral** rather than **pointwise–sectional** language.  
2. **Connection to known rigidity theory:** Known theorems (e.g., Gursky, Hebey–Vaugon, Chang–Gursky–Yang) show partial results where small \(L^2\) or \(L^{n/2}\) Weyl energy under strong curvature restrictions yields local symmetry, but no general threshold theorem is known for the **positively curved Einstein** case. This conjecture refines and unifies those scattered phenomena.  
3. **Interdisciplinary depth:** A solution likely requires techniques from **Ricci flow stability**, **Bochner–Weitzenböck identities for the Weyl tensor**, and **Sobolev–analytic curvature estimates**, linking geometric analysis and algebraic topology of curvature operators.  
4. **Extreme simplicity of form:** Its statement is short, geometrically transparent, and scale–invariant, yet its resolution would reveal profound structure about how conformal geometry controls global topological classification in positive curvature.

Hence, this problem exemplifies a “narrow entrance, profound depth” conjecture—simple to state, technically formidable, and conceptually unifying, making it a truly “good” research-level problem.